\scanpage{100}
\mspage{49}
la ``\uline{découpe}'' de $G$ suivant $C$, obtenue en ``coupant'' comme il a été dit \unsure{les} les milieux des arêtes $\in C$ (ce qui revient à ``enlever'' les arêtes, et à remplacer les côtés correspondants par des pts marqués).

Soit un morphisme de normalisation entre m.graphes
\[
(130)\qquad G'\longrightarrow G\deleted{\unsure{'}}.
\]
Si c'est un morphisme de normalisation \uline{strict}, on voit que $G'$ est stable ssi $G$ l'est. (Dans le cas général, on voit que si $S'\unsure{-|\longrightarrow}S$ est surjectif, et si $G'$ est stable, $G$ l'est aussi.)

\deleted{\unsure{Supp.}} Si $X$ est une courbe spéciale (sur une base quelconque $S$), épinglée par $G$, alors on voit (en procédant comme au début de la section) qu'on lui associe canoniquement une courbe spéciale $G'$-épinglée, d'où un foncteur
\[
(\unsure{130|131})\qquad M_G(S)\longrightarrow M_{G'}(S),
\]
\margin{99}
